\chapter{Material Module Semantic Cleanup}
\label{ch:material-semantic-cleanup}

\section{Purpose}

This chapter documents the first low-risk semantic cleanup applied to the
materials module after the architectural analysis of
Chapter~\ref{ch:material-module-analysis}.  The goal of this phase is not to
replace the existing API, but to make the ontology explicit while preserving
the current executable code and tests.

The cleanup follows three rules:
\begin{itemize}
  \item no behavioural change in the constitutive hot path;
  \item no forced migration of existing code;
  \item new names must coexist with legacy names until the next refactor stage.
\end{itemize}

\section{What Was Changed}

\subsection{Constitutive-space aliases}

The file \texttt{MaterialPolicy.hh} now exposes semantically explicit aliases:
\begin{itemize}
  \item \texttt{ContinuumConstitutiveSpace<N>};
  \item \texttt{BeamConstitutiveSpace<N,Dim>};
  \item \texttt{ShellConstitutiveSpace<N,Dim>}.
\end{itemize}

Alongside them, named aliases now exist for the most common spaces:
\begin{itemize}
  \item \texttt{ThreeDimensionalConstitutiveSpace};
  \item \texttt{TimoshenkoBeamConstitutiveSpace3D};
  \item \texttt{MindlinReissnerShellConstitutiveSpace3D};
  \item and their lower-dimensional counterparts.
\end{itemize}

The legacy names \texttt{ThreeDimensionalMaterial},
\texttt{TimoshenkoBeam3D}, etc.\ remain available as compatibility aliases.

\subsection{Stateful constitutive-site aliases}

The file \texttt{LinealElasticMaterial.hh} now exposes:
\begin{itemize}
  \item \texttt{ConstitutiveSite<Relation,StatePolicy>};
  \item \texttt{ElasticConstitutiveSite<Relation>};
  \item \texttt{InelasticConstitutiveSite<Relation>}.
\end{itemize}

These are semantic aliases for the existing \texttt{MaterialInstance} class.
The old name is intentionally preserved, but the new aliases make the object's
meaning explicit: it is already a \emph{stateful constitutive site}, not a
mere material definition.

\subsection{Erased constitutive handles}

The file \texttt{Material.hh} now exposes the following semantic aliases:
\begin{itemize}
  \item \texttt{ConstitutiveHandle<Space>} for the owning erased wrapper;
  \item \texttt{ConstitutiveConstHandleRef<Space>} for the borrowed const view;
  \item \texttt{ConstitutiveHandleRef<Space>} for the borrowed mutable view.
\end{itemize}

This makes the role of \texttt{Material<Policy>} explicit without breaking the
existing API.

\subsection{Constitutive-site accessors}

\texttt{MaterialPoint} and \texttt{MaterialSection} now expose both naming
families:
\begin{itemize}
  \item legacy accessors: \texttt{material\_cref()}, \texttt{material\_ref()};
  \item semantic accessors: \texttt{constitutive\_cref()},
        \texttt{constitutive\_ref()}.
\end{itemize}

Additional site aliases were introduced:
\begin{itemize}
  \item \texttt{ContinuumConstitutiveSite<Space>};
  \item \texttt{SectionConstitutiveSite<Space,PointDim>}.
\end{itemize}

\section{Ownership Semantics Clarified}

One important contradiction was removed from the documentation of
\texttt{MaterialInstance}.  The previous text suggested flyweight-style sharing
as the default interpretation of the relation handle.  That was not true of the
actual copy semantics.

The code already behaved as follows:
\begin{itemize}
  \item ordinary copy construction deep-clones the relation;
  \item explicit sharing occurs only when the user passes an external
        \texttt{std::shared\_ptr<Relation>}.
\end{itemize}

The comments now match that behaviour.  This is important because implicit
sharing between constitutive sites would be semantically dangerous in nonlinear
analysis.

\section{Why This Stage Matters}

The cleanup may look cosmetic at first sight, but it removes a real source of
architectural ambiguity.  Before this change, the same word ``material'' could
refer to all of the following:
\begin{itemize}
  \item a constitutive space;
  \item a physical material definition;
  \item a stateful constitutive site;
  \item an owning erased runtime handle.
\end{itemize}

After this stage, the code base can distinguish these meanings explicitly
without forcing a flag day migration.

\section{What Was Intentionally Not Changed}

This cleanup does \emph{not} yet:
\begin{itemize}
  \item separate constitutive law and constitutive state into distinct runtime
        objects;
  \item remove \texttt{MemoryState} from the default inelastic path;
  \item move \texttt{FiberSection} out of the materials namespace hierarchy;
  \item change the hot-loop data layout of fiber sections or continuum points.
\end{itemize}

Those changes belong to the next refactor stage and would carry more risk.

\section{Immediate Benefits}

The immediate benefits of this stage are:
\begin{itemize}
  \item clearer architectural vocabulary in code and tests;
  \item more honest documentation of ownership and copying;
  \item a smoother path toward the next law/state/algorithm split;
  \item better traceability for future reinforced-concrete, multiscale, and
        fracture work.
\end{itemize}

\section{Next Step}

The next meaningful refactor should preserve these aliases and attack the real
state-management issue:
\begin{center}
\textbf{constitutive law}
\(\oplus\)
\textbf{integration algorithm}
\(\oplus\)
\textbf{constitutive state}
\end{center}

Only after that separation is explicit will it be reasonable to redesign the
storage layout for large-scale reinforced-concrete simulations, nonlinear
dynamics, and continuum cracking or fracture models.
